
\chapter{OMNI, une base de données regroupant les observations du milieu interplanétaire depuis plusieurs décennies}

OMNI\footnote{OMNI n'est pas un acronyme, le nom souhaite simplement souligner la diversité des sources de relevés} est une base de données de la \acrshort{nasa} qui rassemble les mesures du vent solaire et du champ magnétique interplanétaire faites près de la Terre (\cite{SPDF_OMNIWeb}). 
Elle est très utilisée en physique solaire, géophysique et en météorologie de l'espace. 
Cette section présente comment OMNI est construite, à quoi elle sert et quelles données elle contient.

\section{Description}

\paragraph{Objectifs et création}
L'objectif principal d'OMNI est de permettre l'étude des couplages et relations Terre-Soleil, notamment les effets des variations du vent solaire sur la magnétosphère et l'ionosphère (\cite{OMNI_HRO_doc}).
Pour cette raison, OMNI regroupe dans une même base de données des informations sur le vent solaire, les particules énergétiques et plusieurs indices géomagnétiques. 
OMNI a été créée et est maintenue par J.~H.~King et N.~E.~Papitashvili, au centre de données \acrshort{spdf} (\acrlong{spdf}) de la NASA. 
Le relevé de données commence en Novembre 1963, avec le lancement du satellite \acrshort{imp}~1 (\cite{OMNI2_doc}). 
Or, depuis 1963, on compte plusieurs satellites et instruments de relevé qui se sont succédés. 
OMNI assemble donc les mesures de plusieurs missions pour former une série longue et continue (voir \hyperref[biais]{II.1 - Biais et normalisation}). 

\paragraph{Versions et résolutions}
Il existe deux versions de la base : OMNI et OMNI2. 
OMNI est la version d'origine tandis que OMNI2 en est une version améliorée qui utilise une nouvelle méthode de décalage temporel pour \acrshort{isee}~3, Wind et ACE (voir \hyperref[localisation]{II.1 - Localisation}). 
Elle intègre aussi d'autres paramètres absents de la première version.

Les données OMNI sont également disponibles à plusieurs résolutions temporelles.
La haute résolution (\acrshort{hro}) correspond aux relevés de 1 minute et de 5 minutes depuis 1981. 
La basse résolution correspond à des relevés ou moyennes horaires, journalières et sur 27 jours\footnote{correspondant environ à une rotation solaire} depuis 1963. 
Elle possède davantage de paramètres car certains indices sont actualisés toutes les heures uniquement  (voir \hyperref[paramètres]{II.1 - Paramètres}). 

Pour la suite du rapport, sans indication contraire, on supposera que OMNI est la version 2 avec une résolution horaire. 



\paragraph{Satellites et instruments}
La force d'OMNI est qu'elle regroupe plusieurs instruments de différentes missions spatiales. 
Les données proviennent en effet d'une quinzaine de satellites en orbite autour de la Terre, et des satellites \acrshort{isee}~3, Wind et ACE, au point Lagrange $L_1$ (\cite{OMNI2_doc}). 

Les tableaux~\ref{tab:omni_imf} et~\ref{tab:omni_plasma} donnent, pour chaque mission, le responsable scientifique, la période utilisée et la référence. 

% FIGURE 1 (conseillée) : frise chronologique des satellites utilisés dans OMNI
% (IMP 8 : 1973--2000, Geotail : 1995--2006, Wind : 1995--aujourd'hui,
%  ACE : 1998--2024, DSCOVR : 2016--2019 -- dates tirées de la documentation HRO).
% Elle montre bien que la série est un assemblage de plusieurs missions.
% \begin{figure}[h]
%   \centering
%   \includegraphics[width=0.9\textwidth]{figures/omni_spacecraft_timeline.pdf}
%   \caption{Satellites utilisés pour construire OMNI. \textit{(à compléter)}}
%   \label{fig:omni_timeline}
% \end{figure}


\paragraph{Paramètres}
\label{paramètres}
%il faut refaire ce paragraphe
Le tableau  résume les principaux paramètres des versions 1 min et 5 min. 
Les composantes vectorielles sont données dans les repères GSE et GSM. 
La version horaire contient les mêmes grandes familles de paramètres ainsi que des données utiles sur le long terme : le rapport entre les densités des particules alpha et des protons, des flux de particules énergétiques, le nombre de taches solaires et les indices Kp et Dst (\cite{OMNI2_doc}). 
%textit{(Les indices ap, F10.7 et le nombre de taches solaires sont listés pour les moyennes journalières et sur 27 jours ; à confirmer pour la version horaire.)} 
La version 5 minutes ajoute des flux de protons de plus de 10, 30 et 60~MeV mesurés par les satellites GOES, disponibles seulement en 5 minutes. 
Les indices AE, AL, AU et SYM/H sont calculés à une minute par le centre de données de Kyoto (\cite{OMNI_HRO_doc}).




\paragraph{Localisation}
\label{localisation}
Les satellites n'étant pas tous à la même position spatiale, les données sont donc ramenées à un même point de référence, le nez du bow shock terrestre. 
C'est le point de l'onde de choc situé face au Soleil, juste avant la magnétosphère. 
Pour cela, chaque mesure est décalée dans le temps. 
Les auteurs supposent que le vent solaire transporte des fronts plans et que les mêmes valeurs seront observées plus tard au nez du bow shock (voir \S 12. \cite{OMNI_LRO_doc}). 
\textit{des papiers ont critiqués cette hypothèse linéaire en comparant à des données d'autres satellites geo} \\
La date associée à chaque valeur est donc la date d'arrivée au nez du bow shock, et non la date de mesure. 
Pour les données horaires, le décalage est fait avec des mesures de résolution plus fine (1 à 5 minutes), puis la moyenne horaire est calculée à l'heure terrestre (\cite{OMNI_HRO_doc, OMNI2_doc}).


\paragraph{Biais et normalisation}
\label{biais}
Les instruments des différents satellites ne donnent pas exactement les mêmes valeurs et peuvent introduire des biais dans la continuité.
Pour éviter cela, les chercheurs ont mis en place une normalisation de certaines données (voir \S 13. \cite{OMNI_LRO_doc}). 

Cette normalisation est également accompagnée d'une priorité quand deux satellites fournissent la même information sur la même période (voir \S 8. \cite{OMNI_LRO_doc}). 

\vspace{1cm}

Avec la rapide description de la base de données, on peut se rendre compte que les données accessibles sont déjà pré-traitées. 
Au moment de les télécharger, tous les calculs de décalage temporel, normalisation, priorité, moyenne sont déjà effectués. 
Il ne nous est donc pas nécessaire de devoir le faire et ce n'est d'ailleurs pas l'objet du projet. 
Cependant, il est important d'être avertis et conscients de ces opérations pour de futures hypothèses.  


\section{Données manquantes}

Malgré les efforts des auteurs de la base de données pour prétraiter les données, un problème majeur persiste : les données manquantes. 



